% Generated draft: requires mathematical and source review.
\documentclass[11pt]{article}
\usepackage[margin=25mm]{geometry}
\usepackage{fontspec}
\setmainfont{Latin Modern Roman}
\usepackage{amsmath,amssymb,mathtools,unicode-math}
\setmathfont{Latin Modern Math}
\usepackage{enumitem,needspace,xurl,hyperref}
\setlength{\parindent}{0pt}
\setlength{\parskip}{5pt}
\setlength{\emergencystretch}{3em}
\newcommand{\N}{\mathbb{N}}
\newcommand{\Z}{\mathbb{Z}}
\newcommand{\Q}{\mathbb{Q}}
\newcommand{\R}{\mathbb{R}}
\newcommand{\C}{\mathbb{C}}
\newcommand{\F}{\mathbb{F}}
\newcommand{\A}{\mathbb{A}}
\newcommand{\PP}{\mathbb{P}}
\newcommand{\E}{\mathbb{E}}
\newcommand{\eps}{\varepsilon}
\newcommand{\dd}{\,\mathrm{d}}
\begin{document}
\section*{UnsolvedMath ID: 800811}
\textbf{Source problem code:} catalogue entry
\subsection{Definitions and mathematical statement}
Let \(M\) denote the spatial manifold of a fluid motion, let \(I\subseteq\R\) be its time interval, and let \(u(t,\cdot)\) denote its velocity field. Whenever the associated particle flow is defined, write \(\eta(t,a)\) for the position at time \(t\) of the particle initially at \(a\in M\), so that, for a chosen initial time \(t_*\in I\),
\[
\partial_t\eta(t,a)=u\bigl(t,\eta(t,a)\bigr),\qquad \eta(t_*,a)=a.
\]
These symbols describe the question; they do not specify a fluid equation or an admissible solution class, neither of which is given in the supplied record.

In a specified real-analytic spatial structure, a trajectory is real analytic in time at an interior time \(t_0\in I\) if, in a real-analytic coordinate chart \(\phi\) about \(\eta(t_0,a)\), there exist \(r>0\) and vectors \(c_k\in\R^{\dim M}\) such that \((t_0-r,t_0+r)\subseteq I\), the trajectory stays in the chart on this interval, and
\[
\phi\bigl(\eta(t,a)\bigr)=\sum_{k=0}^{\infty}c_k(t-t_0)^k
\qquad (|t-t_0|<r),
\]
with a convergent power series. For a domain or manifold with boundary, this definition requires an appropriate specified analytic coordinate structure or a specified analytic ambient realization. Arbitrary smooth coordinate changes need not preserve this property.

\textbf{Source question (Problem 10).} Are particle trajectories real analytic in time in every finite spatial dimension? For fluids on manifolds with boundary, how does this property depend on whether the boundary is real analytic?

The intended quantifiers would range over spatial dimensions, admissible fluid motions, particle labels, and times during the motion's existence. However, the supplied statement does not identify the governing equations, the regularity and existence assumptions, the geometric hypotheses, or the boundary conditions that determine these ranges. It also does not say whether boundary particles are included. Thus it does not support an unconditional assertion that every velocity field or every fluid motion has analytic trajectories. For a nonanalytic boundary, the intended meaning of trajectory analyticity must also be specified, for example through an analytic ambient manifold. The question is retained with these unresolved qualifications; it is not a fully specified conjecture.
\subsection{Short English statement}
The source asks whether fluid particles move along trajectories that are real analytic as functions of time in every spatial dimension, and how boundary analyticity affects this property. The fluid equations, solution regularity, geometric setting, and boundary conditions are unspecified, so the draft records the question and its missing hypotheses without selecting a particular fluid model.
\subsection{Sources}
\begin{itemize}
\item Supplied local UnsolvedMath registry record, permanent ID \(800811\), exact source code \texttt{catalogue entry}, titled \emph{Geometric Hydrodynamics in Open Problems --- Problem 10}. The record attributes the problem to Khesin, Misiolek, and Modin, with proposed year 2022.
\item Khesin, Misiolek, and Modin, \emph{Geometric Hydrodynamics in Open Problems} (2022), Problem 10. Supplied source URL: \url{https://arxiv.org/abs/2205.01143}. The record reports access on 2026-10-16 and extraction by \texttt{arxiv\_latex}; these are supplied provenance statements, not an independent source visit or verification.
\item The local record labels the problem \texttt{open}. This is a supplied registry status, not an independently verified assessment of the literature. The supplied upstream revision is \texttt{d371ee917ac9647fbb06cd6fc529deec6f3dda82}.
\end{itemize}
\end{document}
